\section{Relational witnesses and Cohen-model separations}
\label{sec:p06-cardinal-characteristics-separations}

An absolute inequality is established by a transformation of witnesses.
A strict consistency separation needs a model calculation: both an upper
bound supplied by a family and a lower bound ruling out smaller families.
The relational presentation below is that of
\cite[\S\S2, 4, arXiv version 1]{BlassReductions1994V1}.
His Cohen example in Section 2 separates splitting from domination.
The subsequent forcing calculation gives the supporting preservation
arguments and derives the complete classical diagram values for a stated
family of Cohen extensions. It uses the binary Cohen poset already
specified in Table~\ref{tab:p06-forcing-catalog}; it is not a claim that
all assignments compatible with the diagram have thereby been realized.
The changes of underlying real space are also discussed in the separately
dated \cite[pp.~3--4, author draft of 24 November 2003]{BlassCombinatorialAuthor2003}.

\subsection{The relation, its dual, and the direction of a morphism}

\begin{hbdefinition}{def:p06-cc-relational-norm}{Blass's relational norm and dual}
A relation datum is a triple \(\mathbf A=(A_-,A_+,A)\): the third component
\(A\) is a binary relation of type \(A_-\times A_+\), whereas the
triple itself is used as the argument of the norm. Assume that both sets
are nonempty, that every \(x\in A_-\) is related to some \(a\in A_+\),
and that no one \(a\in A_+\) is related to every \(x\in A_-\). Put
\[
\begin{gathered}
 \|\mathbf A\|=\min\{|Z|:Z\subseteq A_+,
                  (\forall x\in A_-)(\exists z\in Z)A(x,z)\},
 \\
 \mathbf A^\perp=(A_+,A_-,\{(a,x):\neg A(x,a)\}).
\end{gathered}
\]
These assumptions ensure that both norms exist and are at least two.
For \(\mathbf W=(\omega^\omega,\omega^\omega,\leq^*)\),
\(\|\mathbf W\|=\mathfrak d\) and \(\|\mathbf W^\perp\|=\mathfrak b\).
Blass writes \(\mathcal P_\infty(\omega)=[\omega]^\omega\) for the
infinite subsets of \(\omega\). For the datum whose negative domain is this set, whose
positive domain is \(\mathcal P(\omega)\), and whose relation says
``is split by,'' its norm is \(\mathfrak s\) and its dual norm is
\(\mathfrak r\). The four cardinals retain
\cref{def:p06-cc-bdsr}'s conventions.
\end{hbdefinition}

\begin{hbdefinition}{def:p06-cc-blass-morphism}{The source's morphism direction}
Let \(\mathbf B=(B_-,B_+,B)\) have the same standing conditions.
A morphism \(\xi:\mathbf A\longrightarrow\mathbf B\) consists of
\[
 \xi_-:B_-\longrightarrow A_-,\qquad
 \xi_+:A_+\longrightarrow B_+
\]
with the implication
\[
 A(\xi_-(b),a)\ \Longrightarrow\ B(b,\xi_+(a)).
\]
Thus its arrow goes in the direction of this implication, and reverses
the induced norm comparison. This is Blass's stated convention, not
an identification with the oppositely directed Vojt\'a\v{s} convention.
A Borel morphism additionally requires Borel domains, relations and
component maps in their stated coding spaces; an arbitrary set map
does not acquire that stronger property from its norm inequality.
\end{hbdefinition}

\begin{hbtheorem}{thm:p06-cc-morphism-duality}{proved=annex:p06.ccsep.morphism-duality}{Witness transfer and reversal under duality}
A morphism \(\mathbf A\to\mathbf B\) yields
\[
 \|\mathbf B\|\leq\|\mathbf A\|,\qquad \|\mathbf A^\perp\|\leq\|\mathbf B^\perp\|.
\]
Interchanging \(\xi_-\) and \(\xi_+\) gives a morphism
\(\mathbf B^\perp\to \mathbf A^\perp\). In particular the interval construction
in \cref{thm:p06-cc-combinatorial} witnesses both
\(\mathfrak s\leq\mathfrak d\) and \(\mathfrak b\leq\mathfrak r\),
not merely the first numerical comparison.
\end{hbtheorem}

\begin{hbdefinition}{def:p06-cc-sequential}{Sequential composition}
For two standing relation data, retain the source's function-space
composition
\[
 \mathbf A;\mathbf B=\bigl(A_-\times{}^{A_+}B_-,\ A_+\times B_+,\ R\bigr),
\]
where \({}^{A_+}B_-\) is the set of all functions \(A_+\to B_-\), and
\[
 R((x,\rho),(a,b))\quad\Longleftrightarrow\quad
       A(x,a)\ \text{and}\ B(\rho(a),b).
\]
The second task therefore receives the answer to the first task as an
input. This higher-type function space is not, in general, a standard
Borel space of real codes. For a morphism \(\mathbf A;\mathbf B\to\mathbf C\),
where \(\mathbf C=(C_-,C_+,C)\) is another such datum, the source's
Borel convention tests the uncurried map \(C_-\times A_+\to B_-\),
as well as the other component maps; it does not pronounce every
function \(A_+\to B_-\) Borel.
\end{hbdefinition}

\begin{hbtheorem}{thm:p06-cc-sequential-norm}{proved=annex:p06.ccsep.sequential-norm}{Product and minimum norms}
For the standing data,
\[
 \|\mathbf A;\mathbf B\|=\|\mathbf A\|\cdot\|\mathbf B\|,\qquad
 \|(\mathbf A;\mathbf B)^\perp\|=\min\{\|\mathbf A^\perp\|,\|\mathbf B^\perp\|\}.
\]
The first expression is cardinal multiplication. For infinite norms it
equals their maximum; finite norms must not be substituted into that
infinite-cardinal simplification.
\end{hbtheorem}

\subsection{Cantor space, real-line ideals, and category witnesses}

\begin{hbconvention}{def:p06-ccsep-cantor}{The domains of the ideal symbols}
Let \(2^\omega\) have its product topology and fair-coin product measure
\(\mu\); a cylinder \([s]=\{x:s\subseteq x\}\),
\(s\in2^{<\omega}\), has measure \(2^{-|s|}\).
Write \(\mathcal B\) for its ideal of meager sets, following
\cite[\S3]{BlassReductions1994V1}.
Here \(\mathcal N\) denotes the ideal of arbitrary subsets of Cantor
space of product outer measure zero, as in
\cite[\S\S2--3]{BartoszynskiInvariants1999V1}. Its domain in this
subsection and the forcing construction is Cantor space, not the real-line
domain of \cref{def:p06-cc-standing}. The bridge theorem below justifies
the comparison of the numerical invariants; it does not identify the
two ideals as sets.
All four ideal invariants use \cref{def:p06-cc-ideal-four}.
In particular arbitrary witnesses need not be Borel or measurable.
\end{hbconvention}

\begin{hbtheorem}{thm:p06-ccsep-domain-bridge}{proved=annex:p06.ccsep.domain-bridge}{Changing the underlying real space}
Each of \(\operatorname{add},\operatorname{cov},\operatorname{non},
\operatorname{cof}\) has the same value for the Cantor null ideal
and the real-line null ideal, and the same value for \(\mathcal B\)
and the real-line meager ideal \(\mathcal M\).
The measure bridge is binary expansion
\(\pi(x)=\sum_{n<\omega}x(n)2^{-n-1}\); the category bridge first
removes the countable exceptional expansions and dyadic points.
\end{hbtheorem}

\begin{hbtheorem}{thm:p06-ccsep-splitting-category}{proved=annex:p06.ccsep.splitting-category}{Splitting and reaping through category}
Every nonmeager subset of \(2^\omega\), viewed as a family of
subsets of \(\omega\), is a splitting family. Consequently
\[
 \mathfrak s\leq\operatorname{non}(\mathcal B),\qquad
 \operatorname{cov}(\mathcal B)\leq\mathfrak r.
\]
These statements quantify over arbitrary infinite sets being split,
including sets that did not belong to a forcing ground.
\end{hbtheorem}

\subsection{The source's chopped-real presentation of category}

\begin{hbdefinition}{def:p06-ccsep-chopped}{Matching and engulfing chopped reals}
Following \cite[\S3]{BlassReductions1994V1}, a chopped real is
\((x,\Pi)\), where \(x\in2^\omega\) and \(\Pi\) partitions
\(\omega\) into consecutive nonempty finite intervals
\([a_n,a_{n+1})\), with \(a_0=0\). Write \(\mathrm{CR}\) for
the set of chopped reals. The set \(\operatorname{Match}(x,\Pi)\)
consists of the \(y\) agreeing with \(x\) on infinitely many
whole intervals of \(\Pi\).
The chopped real \((x,\Pi)\) engulfs \((x',\Pi')\) if all but
finitely many intervals \(I\) of \(\Pi\) contain an interval
\(J\) of \(\Pi'\) on which \(x\) and \(x'\) agree.
This is not a requirement that one partition refine the other everywhere.
\end{hbdefinition}

\begin{hblemma}{thm:p06-ccsep-chopped-base}{proved=annex:p06.ccsep.chopped-base}{A basis of matching sets}
Every \(\operatorname{Match}(x,\Pi)\) is a dense \(G_\delta\).
Every dense \(G_\delta\) subset of Cantor space contains such a
matching set. Consequently every meager set is contained in
\(2^\omega\setminus\operatorname{Match}(x,\Pi)\) for some chopped real.
\end{hblemma}

\begin{hblemma}{thm:p06-ccsep-engulfing}{proved=annex:p06.ccsep.engulfing}{Exact containment criterion}
\[
 \operatorname{Match}(x,\Pi)\subseteq
 \operatorname{Match}(x',\Pi')
 \quad\Longleftrightarrow\quad
 (x,\Pi)\text{ engulfs }(x',\Pi').
\]
\end{hblemma}

\begin{hbdefinition}{def:p06-ccsep-category-data}{The three source relation data}
Retain the source names
\[
\begin{gathered}
 \mathbf U=(\mathrm{CR},2^\omega,\text{matches}^{\smile}),\quad
 \mathbf V=(\mathrm{CR},\mathrm{CR},\text{is engulfed by}),\\
 \mathbf W=(\omega^\omega,\omega^\omega,\leq^*).
\end{gathered}
\]
Thus \(U((x,\Pi),y)\) means \(y\in\operatorname{Match}(x,\Pi)\),
where the source's superscript \(\smile\) takes the converse relation,
and \(V((x,\Pi),(z,\Theta))\) means that \((z,\Theta)\)
engulfs \((x,\Pi)\). Each triple and its binary relation have
the separate argument types just stated. The letter \(V\) here
is not a set-theoretic universe.
\end{hbdefinition}

\begin{hbtheorem}{thm:p06-ccsep-category-norms}{proved=annex:p06.ccsep.category-norms}{The four category cardinals as norms}
\[
 \|\mathbf U\|=\operatorname{non}(\mathcal B),\quad
 \|\mathbf U^\perp\|=\operatorname{cov}(\mathcal B),\quad
 \|\mathbf V\|=\operatorname{cof}(\mathcal B),\quad
 \|\mathbf V^\perp\|=\operatorname{add}(\mathcal B).
\]
\end{hbtheorem}

\begin{hbtheorem}{thm:p06-ccsep-category-sequential}{proved=annex:p06.ccsep.category-sequential}{The Borel sequential morphism for category}
There is a morphism \(\mathbf U;\mathbf W\to\mathbf V\), with the uncurried Borel
standing specified in \cref{def:p06-cc-sequential}. It gives
\[
 \operatorname{add}(\mathcal B)\geq
 \min\{\operatorname{cov}(\mathcal B),\mathfrak b\},\qquad
 \operatorname{cof}(\mathcal B)\leq
 \max\{\operatorname{non}(\mathcal B),\mathfrak d\}.
\]
Combined with the opposite bounds and the domain bridge, these are
the two category identities in their chopped-real presentation.
The second-stage function depends on the first-stage answer \(y\),
not solely on the initial chopped real; omitting that input would
not describe this witness construction.
\end{hbtheorem}

\subsection{The forcing construction and its preservation argument}

\begin{hbconvention}{def:p06-ccsep-ground}{Ground hypotheses for the value calculation}
Use the supplied transitive ground \(M\) and generic-filter semantics
of \cref{def:p06-ftr-ground}. Assume internally that
\[
 M\models\mathsf{ZFC}+\mathsf{CH},\qquad
 \kappa>\aleph_1\ \text{is regular},\qquad
 \kappa^{\aleph_0}=\kappa.
\]
Let
\(P_\kappa=\operatorname{Fn}(\kappa\times\omega,2,{<}\omega)^M\),
ordered by reverse inclusion, and let \(G\) be supplied generic over
\(M\). Put \(c_\alpha(n)=(\bigcup G)(\alpha,n)\).
The union is total by the coordinate dense sets.
For \(S\subseteq\kappa\) in \(M\), set
\(P_S=\operatorname{Fn}(S\times\omega,2,{<}\omega)^M\), and
\(G_S=\{p\mathbin{\upharpoonright}(S\times\omega):p\in G\}\).
The restricted conditions are regarded as conditions in \(P_S\).
All bounds and regularity hypotheses on \(\kappa\) are ground-model
hypotheses; preservation below explains their subsequent use in \(M[G]\).
No transitive ground or external generic is inferred merely from a
consistency statement.
\end{hbconvention}

\begin{hblemma}{thm:p06-ccsep-support}{proved=annex:p06.ccsep.support}{Countable names and fresh coordinates}
The forcing \(P_\kappa\) is ccc, preserves ground cardinals and
cofinalities, and forces \(\mathfrak c=\kappa\).
Every real or countable sequence of real codes in \(M[G]\) belongs
to \(M[G_S]\) for some countable \(S\subseteq\kappa\) in \(M\).
More generally every family of fewer than \(\kappa\) such objects
is contained in one \(M[G_S]\) with \(S\in M\) and
\(|S|^M<\kappa\). If \(\alpha\notin S\), then
\(c_\alpha\) is Cohen generic over \(M[G_S]\).
The set \(S\) need not be an initial segment of \(\kappa\).
\end{hblemma}

\begin{hblemma}{thm:p06-ccsep-category-preservation}{proved=annex:p06.ccsep.category-preservation}{Countable forcing preserves old nonmeager sets}
If a countable forcing \(P\in M\) is used and \(A\subseteq2^\omega\)
is nonmeager in \(M\), then \(A\) remains nonmeager in the extension.
For \(P_\kappa\), the old Cantor reals \((2^\omega)^M\) remain
nonmeager. Thus, under the standing \(\mathsf{CH}\),
\(\operatorname{non}(\mathcal B)=\aleph_1\).
The proof handles the newly coded meager cover; it does not assume
that every nonmeager set is preserved by every ccc forcing.
\end{hblemma}

\begin{hblemma}{thm:p06-ccsep-fresh-obstructions}{proved=annex:p06.ccsep.fresh-obstructions}{One fresh Cohen coordinate defeats three small families}
In the standing extension,
\[
 \operatorname{cov}(\mathcal B)=\mathfrak d=\mathfrak r=\kappa.
\]
A fresh coordinate avoids a meager cover coded in the smaller extension,
supplies an unbounded function against its dominating family, and splits
every member of its family of infinite subsets of \(\omega\).
These are three separate density arguments.
The function supplied by the second argument is \(f_c(n)\), the
position of the \(n\)-th one in the fresh binary real \(c\), counted
from zero; genericity ensures that every such position exists.
\end{hblemma}

\begin{hbdefinition}{def:p06-ccsep-cohen-null}{The null set coded by a Cohen real}
Put \(a_0=0\), \(a_{n+1}=a_n+n+1\) and
\(I_n=[a_n,a_{n+1})\). For \(c\in2^\omega\), let
\[
 N_c=\bigcap_{m<\omega}\bigcup_{n\geq m}
             \{x\in2^\omega:x\mathbin{\upharpoonright}I_n
                          =c\mathbin{\upharpoonright}I_n\}.
\]
Each event in the union has measure \(2^{-n-1}\).
The subscripted set \(N_c\) is a subset of Cantor space, not a model
or an ideal. Its definition uses infinitely many matches, unlike the
eventual localization relation of \cref{def:p06-cc-slaloms}.
\end{hbdefinition}

\begin{hblemma}{thm:p06-ccsep-null-values}{proved=annex:p06.ccsep.null-values}{The null covering and uniformity values}
For every \(c\), the set \(N_c\) is null. If \(c\) is Cohen over
\(M'\), then \((2^\omega)^{M'}\subseteq N_c\).
In the standing \(P_\kappa\) extension,
\[
 \operatorname{cov}(\mathcal N)=\aleph_1,
 \qquad \operatorname{non}(\mathcal N)=\kappa.
\]
In particular the old real set can be null and nonmeager simultaneously;
there is no contradiction because the ideals are different.
\end{hblemma}

\begin{hbtheorem}{thm:p06-ccsep-cohen-values}{proved=annex:p06.ccsep.cohen-values}{The complete classical diagram calculation in this Cohen model}
Transporting the ideal values to the real line, the standing extension
satisfies
\[
\begin{aligned}
 \operatorname{add}(\mathcal N)
 &=\operatorname{cov}(\mathcal N)
 =\operatorname{add}(\mathcal M)
 =\operatorname{non}(\mathcal M)
 =\mathfrak b=\mathfrak s=\aleph_1,\\
 \operatorname{non}(\mathcal N)
 &=\operatorname{cof}(\mathcal N)
 =\operatorname{cov}(\mathcal M)
 =\operatorname{cof}(\mathcal M)
 =\mathfrak d=\mathfrak r=\mathfrak c=\kappa.
\end{aligned}
\]
For a ground satisfying \(\mathsf{GCH}\), one may take
\(\kappa=\aleph_2\). This gives, among other separations,
\(\mathfrak s<\mathfrak d\), \(\mathfrak b<\mathfrak r\),
\(\operatorname{cov}(\mathcal N)<\operatorname{cov}(\mathcal M)\),
and \(\operatorname{non}(\mathcal M)<\operatorname{non}(\mathcal N)\).
The construction supplies an extension of the stated ground; it does
not assert that such a ground is a theorem of \(\mathsf{ZFC}\).
\end{hbtheorem}

\begin{hbwarning}{def:p06-ccsep-separation-boundary}{Which comparisons have been separated}
The displayed values realize one pattern, not every permitted pattern.
They do not separate \(\mathfrak b\) from \(\mathfrak s\), nor
\(\mathfrak d\) from \(\mathfrak r\), in this model. The equations
for category remain the absolute identities of
\cref{thm:p06-cc-category-identities}; their failure is not an available
forcing target. Preservation of a measure-positive set, preservation
of a dominating family, and preservation of properness are different
properties and require their own arguments for another forcing.
\end{hbwarning}
